\section*{Chapter \ref{ch:iv:phone-and-technical-screens} --- Phone and Technical Screens}

\begin{solution}{iq:iv:phone-and-technical-screens:1}
The larger of two dice is at least the first die and strictly larger whenever the second is higher, so its
expectation exceeds the first die's 3.5: a range check. $\P(\max \le k) = (k/6)^2$, so $\P(\max = k) = (2k -
1)/36$ and $\E[\max] = \sum_k k(2k-1)/36 = 161/36 \approx 4.47$.
\iqlookfor{a one-sentence argument that rejects the answer, then the distribution of the maximum.}
\end{solution}

\begin{solution}{iq:iv:phone-and-technical-screens:2}
A probability cannot exceed one: adding the ten probabilities counts the days on which several alerts fire
more than once. The complement gives $1 - 0.85^{10} \approx 0.803$.
\iqlookfor{the range check and the complement.}
\end{solution}

\begin{solution}{iq:iv:phone-and-technical-screens:3}
With four prices and $k = 2$ the loop runs for $i = 0, 1$ and returns $[10.5, 11.5]$; the windows are
$[10, 11], [11, 12], [12, 13]$, so the answer should be $[10.5, 11.5, 12.5]$. There are $n - k + 1$ windows
and the loop stops one early (with $n = k$ it returns nothing). Fix the bound to \texttt{range(len(prices) - k
+ 1)}, and for $O(n)$ keep a running sum: add the entering price, subtract the leaving one. Guard $k \le 0$
and $k > n$. The chapter's code tests the fix against a direct computation on 200 random cases.
\iqlookfor{running a small case by hand, the count of windows, and the running-sum improvement.}
\end{solution}

\begin{solution}{iq:iv:phone-and-technical-screens:4}
A portfolio's volatility cannot exceed the weighted average of its components' volatilities, reached only at
correlation 1; that average is 20\%, so 25\% is impossible. The variance is $0.25 \times 0.04 + 0.25 \times
0.04 + 2 \times 0.25 \times 0.5 \times 0.04 = 0.03$, so the volatility is $\sqrt{0.03} \approx 17.3\%$.
\iqlookfor{the bound from the triangle inequality and the two-asset variance formula.}
\end{solution}

\begin{solution}{iq:iv:phone-and-technical-screens:5}
Small cases: $n = 1$ gives 1, $n = 2$ gives 2, $n = 3$ gives 3, where the conjecture says 4. It is refuted at
$n = 3$. The last column is either one vertical domino, leaving a $2 \times (n-1)$ board, or two horizontal
dominoes, leaving $2 \times (n-2)$: $T_n = T_{n-1} + T_{n-2}$ with $T_1 = 1$, $T_2 = 2$, the Fibonacci numbers
1, 2, 3, 5, 8, 13.
\iqlookfor{testing a conjecture on small cases before stating it, and the recursion on the last column.}
\end{solution}

\begin{solution}{iq:iv:phone-and-technical-screens:6}
Say where you are and what blocks you, specifically: ``I have the SDE for the discounted price and I need its
expectation; I am not sure whether the Itô term vanishes here. Would it help to take the expectation before
applying the formula?'' Then use the hint and finish aloud. Avoid silence, avoid restarting from scratch
without saying why, and avoid bluffing a step you cannot justify. If the time runs out, state the answer you
would expect and how you would check it: an interviewer can score a correct plan.
\iqlookfor{a specific request for help, and continued structured progress.}
\end{solution}

\begin{solution}{iq:iv:phone-and-technical-screens:7}
The error is in the re-roll branch: when you re-roll you get a fresh die worth 3.5 on average, not 4. Re-roll
when the first roll is 1, 2 or 3 (below 3.5): the value is $\tfrac12 \times 5 + \tfrac12 \times 3.5 = 4.25$,
which is $17/4$. With three rolls, the last two are worth 4.25, so keep a first roll of 5 or 6 and re-roll
otherwise: $\tfrac13 \times 5.5 + \tfrac23 \times 4.25 = 14/3 \approx 4.67$.
\iqlookfor{backward induction on the continuation value, and spotting the double-counted average.}
\end{solution}

\begin{solution}{iq:iv:phone-and-technical-screens:8}
Three checks. Bounds: at least 6 rolls, and the last new face alone takes 6 rolls on average, so the answer is
well above 6. Structure: after $k$ faces are seen, a new one appears with probability $(6 - k)/6$, so the wait
is geometric with mean $6/(6-k)$; summing gives $6(1 + \tfrac12 + \dots + \tfrac16) = 14.7$. Small case: for a
coin ($n = 2$) the formula gives $2(1 + \tfrac12) = 3$, which is right (one toss, then a wait of mean 2 for
the other side). The chapter's simulation over ten seeds agrees.
\iqlookfor{bounds, the geometric decomposition and a small-case check, all stated in seconds.}
\end{solution}
